\begin{textsample}{POS Dim 2 – ac – Score 25.00 – ac/ac\_em\_12.txt}  \label{ex:f2_pos_166}
9.3.
Ensino em tempo parcial A \textbf{Reforma} do Ensino Médio surgiu a partir da necessidade de ter uma escola mais conectada com as demandas dos estudantes da geração atual.
O modelo \textbf{instituído} pela Lei nº13.415/2017 altera as \textbf{Diretrizes} e Bases da Educação Nacional e se destaca pela \textbf{carga} \textbf{horária} do Novo Ensino Médio, que trouxe flexibilidade para a organização curricular.
Nesse sentido, a \textbf{oferta} de ensino em tempo parcial se antes a antiga \textbf{carga} \textbf{horária} do ensino médio era de 2.400 horas, agora será de 3.000 horas, sendo que 1.800 horas serão \textbf{destinadas} para as aprendizagens comuns e obrigatórias \textbf{previstas} pela BNCC, e as 1.200 horas são voltadas para o \textbf{itinerário} formativo, sendo \textbf{ofertada} nas escolas de tempo parcial em dois turnos para diferentes públicos.
No estado do Acre, segundo dados do senso escolar existe 236 ( \textbf{duzentos} e trinta e seis ) escola com \textbf{oferta} de ensino em tempo parcial, além dessas tem as escolas de ensino militar que é uma unidade escolar estadual integrada à estrutura organizacional da \textbf{instituição} Corpo de Bombeiros Militar do Estado do Acre – CBMAC, conforme Art. 1º da Lei Estadual 3.362/2017.
Seu funcionamento se dá em regime de colaboração com a Secretaria de Educação e Esporte – SEE, conforme Art. 1 §1º da referida lei.
E, as escolas cívico militares, criada pelo Decreto n{\EmojiFont °} 10.356, de 14 de dezembro de 2016, mas seu funcionamento só ocorreu em 2020, na \textbf{gestão} do governador Gladson Cameli, por meio da \textbf{Portaria} n{\EmojiFont °} 345.
A \textbf{instituição} educacional tem como parceiros o Corpo de Bombeiros Militar do Acre e a Polícia Militar do Estado do Acre, \textbf{atendendo} estudantes de Rio Branco e \textbf{municípios}. 9.4.
Ensino em tempo integral A Educação em Tempo Integral é uma idealização do \textbf{governo} \textbf{federal} que garante o desenvolvimento dos estudantes em todas as dimensões formativas : intelectual, física, afetiva, social e cultural, a partir da ampliação da jornada escolar, com a condição de que as aprendizagens ocorram de maneira criativa e articulada, a partir de um Projeto Político Pedagógico e Plano de Ação elaborados coletivamente, no qual os objetivos educacionais sejam compartilhados por todos que fazem parte da comunidade escolar, de modo que se sintam corresponsáveis pelas metas e ações planejadas.
No ano de 2017, foram implantadas e implementadas, na Rede Pública Estadual de Ensino do Acre, as primeiras Escolas Estaduais de Ensino Médio em Tempo Integral, denominadas “ Escola Jovem ”, vinculadas à Secretaria de Estado de Educação, Cultura e Esportes ( SEE ), objetivando planejar, executar e avaliar ações inovadoras em conteúdo, metodologia e \textbf{gestão}, direcionadas à melhoria da \textbf{oferta} e da qualidade do ensino médio, buscando desenvolver as competências e habilidades exigidas no mundo contemporâneo.
Essa proposta de educação visa fortalecer o processo formativo, a partir de uma Base Nacional Comum Curricular e uma parte \textbf{flexível}, considerando as \textbf{diretrizes} e parâmetros nacionais e locais, que serão \textbf{ofertadas} em 7 horas diárias, no mínimo, perfazendo uma \textbf{carga} \textbf{horária} anual de 1400 horas de efetiva jornada escolar, sob a mediação dos professores, propondo inovações em conteúdo ( o que ensinar ), método ( como ensinar ), e na \textbf{gestão} dos processos de ensino e de aprendizagem, em diversos espaços sociais e ao longo da vida.
Essas inovações são fundamentadas em princípios educativos que possibilitam as condições necessárias para que os estudantes aprendam a elaborar o seu Projeto de Vida, por meio da ampliação, diversificação e enriquecimento dos conhecimentos e experiências vivenciadas ao longo dos anos do Ensino Fundamental e Ensino Médio, a partir dos seguintes princípios : Protagonismo, Quatro Pilares da Educação, Pedagogia da Presença e Educação Interdimensional, objetivando formar estudantes autônomos, solidários e competentes. 9.5.
Arquitetura curricular Para \textbf{atender} a ampliação do tempo na jornada escolar o módulo/aula passa de 50 ( cinquenta ) para 60 ( \textbf{sessenta} ) minutos, totalizando uma \textbf{carga} \textbf{horária} de 3000 horas para as escolas de ensino médio em tempo parcial diurno e 4.440 horas para as escolas de ensino médio em tempo integral.
Passando a ficar em vigência na Rede Estadual de Educação diante a aprovação do Conselho Estadual de Educação e implementada nas escolas da Rede Estadual de Educação a partir de 2022 de forma gradativa ( 2022 – 1ª \textbf{série} ; 2023 – 2ª \textbf{série} ; 2024 – 3ª \textbf{série} ) para todos os modelos de \textbf{oferta} de ensino médio.
Para a implementação do Novo Ensino Médio é importante destacar que foi elaborada uma arquitetura curricular específica à luz da BNCC, \textbf{atendendo} a parte da Formação Geral Básica e o Itinerário Formativo.
Sendo 1800 horas de Formação Geral Básica para os dois modelos de ensino médio ( Ensino Médio Regular Diurno e Ensino Médio em Tempo Integral ) e 1200 horas de Itinerário Formativo para o Ensino Médio Regular Diurno e 2640 horas para o Ensino Médio em Tempo Integral.
A fim de \textbf{atender} o modelo de Ensino Médio Regular Diurno e Ensino Médio em Tempo Integral, foi preciso desenhar uma proposta de arquitetura curricular específica na Rede Estadual de Educação, considerando o modelo, a \textbf{carga} \textbf{horária}, a realidade local das escolas, o percurso do Itinerário Formativo de \textbf{oferta} Propedêutica ou de Formação \textbf{Técnica} Profissionalizante e ao mesmo tempo, garantir a \textbf{oferta} da Formação Geral Básica comum aos dois modelos.

% matched lemmas: atender, carga, destinar, diretriz, duzentos, federal, flexível, gestão, governo, horário, instituir, instituição, itinerário, município, oferta, ofertar, portaria, prever, reforma, sessenta, série, técnico
\end{textsample}
